\documentclass[11pt]{article}
\usepackage[margin=1in]{geometry}
\usepackage{amsmath, amssymb}
\usepackage{array}

\title{Advanced Algebra II and Precalculus Cheat Sheet}
\author{}
\date{}

\begin{document}
\maketitle

\section*{Laws of Sines and Cosines}

\subsection*{1. Law of Sines}

Use Law of Sines when you know an angle and its opposite side.

\[
\frac{a}{\sin A}
=
\frac{b}{\sin B}
=
\frac{c}{\sin C}
\]

Equivalent form:

\[
\frac{\sin A}{a}
=
\frac{\sin B}{b}
=
\frac{\sin C}{c}
\]

Example:

\[
A=40^\circ,\quad a=10,\quad B=70^\circ
\]

Find \(b\):

\[
\frac{a}{\sin A}=\frac{b}{\sin B}
\]

\[
\frac{10}{\sin40^\circ}
=
\frac{b}{\sin70^\circ}
\]

\[
b=
\frac{10\sin70^\circ}{\sin40^\circ}
\]

\subsection*{2. Law of Cosines}

Use Law of Cosines when you know two sides and the included angle, or all three sides.

\[
c^2=a^2+b^2-2ab\cos C
\]

Example:

\[
a=5,\quad b=7,\quad C=60^\circ
\]

\[
c^2=5^2+7^2-2(5)(7)\cos60^\circ
\]

\[
c^2=25+49-70\cdot\frac{1}{2}
\]

\[
c^2=39
\]

\[
c=\sqrt{39}
\]

\subsection*{3. Solve a Triangle}

To solve a triangle means to find all missing sides and angles.

Always remember:

\[
A+B+C=180^\circ
\]

Checklist:

\begin{enumerate}
  \item If you have an opposite side-angle pair, use Law of Sines.
  \item If you have two sides and the included angle, use Law of Cosines.
  \item If you have all three sides, use Law of Cosines to find an angle.
  \item Use \(A+B+C=180^\circ\) to finish.
\end{enumerate}

\subsection*{4. Area of a Triangle: Sine Formula}

Use this when you know two sides and the included angle.

\[
K=\frac{1}{2}ab\sin C
\]

Example:

\[
a=8,\quad b=10,\quad C=30^\circ
\]

\[
K=\frac{1}{2}(8)(10)\sin30^\circ
\]

\[
K=40\cdot\frac{1}{2}=20
\]

\section*{Ellipses on the Coordinate Plane}

An ellipse is an oval. The center is \((h,k)\).

\subsection*{Horizontal Ellipse}

\[
\frac{(x-h)^2}{a^2}+\frac{(y-k)^2}{b^2}=1
\qquad
a>b
\]

\[
\text{center}=(h,k)
\]

\[
\text{vertices}=(h\pm a,k)
\]

\[
\text{co-vertices}=(h,k\pm b)
\]

\[
c^2=a^2-b^2
\]

\[
\text{foci}=(h\pm c,k)
\]

\subsection*{Vertical Ellipse}

\[
\frac{(x-h)^2}{b^2}+\frac{(y-k)^2}{a^2}=1
\qquad
a>b
\]

\[
\text{center}=(h,k)
\]

\[
\text{vertices}=(h,k\pm a)
\]

\[
\text{co-vertices}=(h\pm b,k)
\]

\[
c^2=a^2-b^2
\]

\[
\text{foci}=(h,k\pm c)
\]

\subsection*{1. Find the Center, Vertices, or Co-Vertices}

Example:

\[
\frac{(x-2)^2}{25}+\frac{(y+1)^2}{9}=1
\]

\[
\text{center}=(2,-1)
\]

Since \(25>9\), the major axis is horizontal.

\[
a^2=25,\quad a=5
\]

\[
b^2=9,\quad b=3
\]

\[
\text{vertices}=(2\pm5,-1)=(-3,-1),(7,-1)
\]

\[
\text{co-vertices}=(2,-1\pm3)=(2,-4),(2,2)
\]

\subsection*{2. Find the Length of the Major or Minor Axis}

\[
\text{major axis length}=2a
\]

\[
\text{minor axis length}=2b
\]

For:

\[
\frac{(x-2)^2}{25}+\frac{(y+1)^2}{9}=1
\]

\[
2a=10,\quad 2b=6
\]

\subsection*{3. Find the Foci}

\[
c^2=a^2-b^2
\]

For:

\[
a^2=25,\quad b^2=9
\]

\[
c^2=25-9=16
\]

\[
c=4
\]

Horizontal ellipse:

\[
\text{foci}=(2\pm4,-1)=(-2,-1),(6,-1)
\]

\subsection*{4. Write Equations of Ellipses from Graphs}

From a graph:

\begin{enumerate}
  \item Find the center \((h,k)\).
  \item Find \(a\), the distance from center to a vertex.
  \item Find \(b\), the distance from center to a co-vertex.
  \item Put the bigger denominator under the direction of the major axis.
\end{enumerate}

Example: center \((1,2)\), horizontal vertex distance \(4\), vertical co-vertex distance \(3\).

\[
\frac{(x-1)^2}{16}+\frac{(y-2)^2}{9}=1
\]

\subsection*{5. Write Equations of Ellipses Using Properties}

Example: center \((-2,5)\), vertices \((-2,9)\) and \((-2,1)\), minor axis length \(6\).

The vertices move vertically, so the major axis is vertical.

\[
a=4,\quad a^2=16
\]

\[
2b=6,\quad b=3,\quad b^2=9
\]

\[
\frac{(x+2)^2}{9}+\frac{(y-5)^2}{16}=1
\]

\subsection*{6. Convert General Form to Standard Form}

General form often looks like:

\[
Ax^2+By^2+Cx+Dy+E=0
\]

Steps:

\begin{enumerate}
  \item Group \(x\)-terms and \(y\)-terms.
  \item Factor coefficients from squared groups.
  \item Complete the square.
  \item Move the constant.
  \item Divide so the right side becomes \(1\).
\end{enumerate}

Example:

\[
4x^2+9y^2-16x+18y-11=0
\]

\[
4(x^2-4x)+9(y^2+2y)=11
\]

\[
4(x^2-4x+4)+9(y^2+2y+1)=11+16+9
\]

\[
4(x-2)^2+9(y+1)^2=36
\]

\[
\frac{(x-2)^2}{9}+\frac{(y+1)^2}{4}=1
\]

\subsection*{7. Find Properties from General Form}

First convert to standard form, then read center, vertices, co-vertices, axes, and foci.

\section*{Hyperbolas on the Coordinate Plane}

A hyperbola has two branches. The center is \((h,k)\).

\subsection*{Horizontal Hyperbola}

\[
\frac{(x-h)^2}{a^2}-\frac{(y-k)^2}{b^2}=1
\]

\[
\text{vertices}=(h\pm a,k)
\]

\[
c^2=a^2+b^2
\]

\[
\text{foci}=(h\pm c,k)
\]

\[
\text{asymptotes: }y-k=\pm\frac{b}{a}(x-h)
\]

\subsection*{Vertical Hyperbola}

\[
\frac{(y-k)^2}{a^2}-\frac{(x-h)^2}{b^2}=1
\]

\[
\text{vertices}=(h,k\pm a)
\]

\[
c^2=a^2+b^2
\]

\[
\text{foci}=(h,k\pm c)
\]

\[
\text{asymptotes: }y-k=\pm\frac{a}{b}(x-h)
\]

\subsection*{1. Find the Center}

Example:

\[
\frac{(x-3)^2}{16}-\frac{(y+2)^2}{9}=1
\]

\[
\text{center}=(3,-2)
\]

\subsection*{2. Find the Vertices}

Since the positive term is the \(x\)-term, the hyperbola opens left and right.

\[
a^2=16,\quad a=4
\]

\[
\text{vertices}=(3\pm4,-2)=(-1,-2),(7,-2)
\]

\subsection*{3. Length of Transverse or Conjugate Axes}

\[
\text{transverse axis length}=2a
\]

\[
\text{conjugate axis length}=2b
\]

For:

\[
\frac{(x-3)^2}{16}-\frac{(y+2)^2}{9}=1
\]

\[
2a=8,\quad 2b=6
\]

\subsection*{4. Equations for the Asymptotes}

Horizontal hyperbola:

\[
y-k=\pm\frac{b}{a}(x-h)
\]

Example:

\[
\frac{(x-3)^2}{16}-\frac{(y+2)^2}{9}=1
\]

\[
h=3,\quad k=-2,\quad a=4,\quad b=3
\]

\[
y+2=\pm\frac{3}{4}(x-3)
\]

\subsection*{5. Find the Foci}

For hyperbolas:

\[
c^2=a^2+b^2
\]

Example:

\[
a^2=16,\quad b^2=9
\]

\[
c^2=16+9=25
\]

\[
c=5
\]

Horizontal hyperbola:

\[
\text{foci}=(3\pm5,-2)=(-2,-2),(8,-2)
\]

\subsection*{6. Write Equations from Graphs}

From a graph:

\begin{enumerate}
  \item Find the center \((h,k)\).
  \item Decide if it opens left/right or up/down.
  \item Find \(a\), the distance from center to vertex.
  \item Find \(b\) using the rectangle/asymptote slope, if given.
  \item Put the opening direction first and positive.
\end{enumerate}

\subsection*{7. Write Equations Using Properties}

Example: center \((0,1)\), opens up/down, vertex distance \(3\), conjugate half-length \(2\).

\[
a=3,\quad b=2
\]

Vertical hyperbola:

\[
\frac{(y-1)^2}{9}-\frac{x^2}{4}=1
\]

\subsection*{8. Convert General Form to Standard Form}

Steps:

\begin{enumerate}
  \item Group \(x\)-terms and \(y\)-terms.
  \item Factor coefficients.
  \item Complete the square.
  \item Move the constant.
  \item Divide so the right side becomes \(1\).
\end{enumerate}

Example:

\[
9x^2-4y^2-54x-16y+29=0
\]

\[
9(x^2-6x)-4(y^2+4y)=-29
\]

\[
9(x^2-6x+9)-4(y^2+4y+4)=-29+81-16
\]

\[
9(x-3)^2-4(y+2)^2=36
\]

\[
\frac{(x-3)^2}{4}-\frac{(y+2)^2}{9}=1
\]

\subsection*{9. Find Properties from General Form}

First convert to standard form, then read the center, opening direction, vertices, axis lengths, asymptotes, and foci.

\section*{Probability}

\subsection*{1. Counting Principle}

If one choice has \(m\) options and another has \(n\) options, total outcomes are:

\[
m\cdot n
\]

Example: \(4\) shirts and \(3\) pants.

\[
4\cdot3=12
\]

\subsection*{2. Combinations and Permutations}

Permutation means order matters.

\[
nP r=\frac{n!}{(n-r)!}
\]

Combination means order does not matter.

\[
nC r=\frac{n!}{r!(n-r)!}
\]

Example:

\[
5P2=\frac{5!}{3!}=5\cdot4=20
\]

\[
5C2=\frac{5!}{2!3!}=10
\]

\subsection*{3. Probabilities Using Combinations and Permutations}

\[
P(\text{event})=
\frac{\text{favorable outcomes}}{\text{total outcomes}}
\]

Example: choose \(2\) students from \(5\). Probability that two specific students are chosen:

\[
\text{total}=5C2=10
\]

\[
\text{favorable}=1
\]

\[
P=\frac{1}{10}
\]

\section*{Binomial Probability Distributions}

Use binomial probability when:

\begin{enumerate}
  \item There are a fixed number of trials, \(n\).
  \item Each trial has success or failure.
  \item Probability of success is constant, \(p\).
  \item Trials are independent.
\end{enumerate}

\subsection*{1. Find Probabilities Using the Binomial Distribution}

\[
P(X=k)=
{n\choose k}p^k(1-p)^{n-k}
\]

Example: \(n=5\), \(p=0.4\), find \(P(X=2)\).

\[
P(X=2)=
{5\choose2}(0.4)^2(0.6)^3
\]

\[
=10(0.16)(0.216)=0.3456
\]

\subsection*{2. Mean, Variance, and Standard Deviation}

For a binomial distribution:

\[
\mu=np
\]

\[
\sigma^2=np(1-p)
\]

\[
\sigma=\sqrt{np(1-p)}
\]

Example: \(n=20\), \(p=0.3\).

\[
\mu=20(0.3)=6
\]

\[
\sigma^2=20(0.3)(0.7)=4.2
\]

\[
\sigma=\sqrt{4.2}
\]

\subsection*{3. Normal Approximation to Binomial}

Use normal approximation when:

\[
np\ge 10
\]

and

\[
n(1-p)\ge 10
\]

Mean and standard deviation:

\[
\mu=np
\]

\[
\sigma=\sqrt{np(1-p)}
\]

Use continuity correction:

\[
P(X\le k)\approx P(Y<k+0.5)
\]

\[
P(X\ge k)\approx P(Y>k-0.5)
\]

\[
P(X=k)\approx P(k-0.5<Y<k+0.5)
\]

\section*{Distributions of Sample Means}

\subsection*{1. Distributions of Sample Means}

If a population has mean \(\mu\) and standard deviation \(\sigma\), then sample means have:

\[
\mu_{\bar{x}}=\mu
\]

\[
\sigma_{\bar{x}}=\frac{\sigma}{\sqrt{n}}
\]

The value \(\sigma_{\bar{x}}\) is called the standard error.

Example:

\[
\mu=80,\quad \sigma=12,\quad n=36
\]

\[
\mu_{\bar{x}}=80
\]

\[
\sigma_{\bar{x}}=\frac{12}{\sqrt{36}}=2
\]

\subsection*{2. The Central Limit Theorem}

The Central Limit Theorem says that for large enough \(n\), the distribution of sample means is approximately normal.

Usually:

\[
n\ge30
\]

is considered large enough.

Use:

\[
z=\frac{\bar{x}-\mu}{\sigma/\sqrt{n}}
\]

Example:

\[
\mu=100,\quad \sigma=15,\quad n=36,\quad \bar{x}=105
\]

\[
z=
\frac{105-100}{15/\sqrt{36}}
\]

\[
=\frac{5}{15/6}
=\frac{5}{2.5}=2
\]

\section*{Quick Formula Sheet}

\[
\frac{a}{\sin A}=\frac{b}{\sin B}=\frac{c}{\sin C}
\]

\[
c^2=a^2+b^2-2ab\cos C
\]

\[
K=\frac{1}{2}ab\sin C
\]

\[
\text{Ellipse: }c^2=a^2-b^2
\]

\[
\text{Hyperbola: }c^2=a^2+b^2
\]

\[
\text{Ellipse major axis}=2a,\qquad \text{minor axis}=2b
\]

\[
\text{Hyperbola transverse axis}=2a,\qquad \text{conjugate axis}=2b
\]

\[
nPr=\frac{n!}{(n-r)!}
\]

\[
nCr=\frac{n!}{r!(n-r)!}
\]

\[
P(X=k)={n\choose k}p^k(1-p)^{n-k}
\]

\[
\mu=np,\qquad \sigma=\sqrt{np(1-p)}
\]

\[
\mu_{\bar{x}}=\mu,\qquad \sigma_{\bar{x}}=\frac{\sigma}{\sqrt{n}}
\]

\[
z=\frac{\bar{x}-\mu}{\sigma/\sqrt{n}}
\]

\end{document}
